% !TeX root = ../papers/nc-p3.tex
\section{十七点支持の曲線状有限窓と接平面の接触次数}
\label{sec:finite-seventeen-site-pencils}
\begin{lemma}\label{lem:embedded-conic-contact-budget}
命題\ref{prop:curvilinear-marked-pencils}の五空間の階数一切断が平面二次曲線、
一つの単純な孤立点と一つの単純な埋込み接点からなる場合を考える。
$J$ の外に十以上の支持があるとする。核側の二位置の各々で、単純印は非単純点のみに
支持される長さ2部分スキームの割線上にある。それが一つの局所因子に集中するには、
その因子は長さ4以上を要する。
\end{lemma}
\begin{proof}
$Z$ 上の独立な相対鉛筆を用いる。孤立点を通る剰余鉛筆は全長でその点を通り、
その固定行因子による平面が捉える印外の長さは6以下である。
外向き接線を持つ剰余鉛筆は全長で固定接平面にある。
この接平面は少なくとも一方の超平面で直線以下を切る（さもなければ外の支持は $2+6$ 以下）。
その直線で数えられる長さは3以下である。

二次曲線の平面 $D$ に入る剰余鉛筆では、局所基底環を $R=k[t]/t^s$ とする。
埋込み接点を避ければ全長で $D$ 内にあり、その長さ $f=s$ を固定左核の直線に数えられる。
接点と別の点を切るなら、接点に延びる単根の法線座標 $z$ は $z^2=0$ を満たす。
$f=\operatorname{ord}z$ は $2f\ge s$ を満たし、長さ $f$ で両根は $D$ にあり、再び核が固定される。

最後に曲線の内側の接線として埋込み接点に集中する剰余鉛筆を扱う。局所模型
$k[[x,z]]/(xz,z^2)$ で行列式の二点因子の代数は $R$ 上 $1,x$ を基底とする。
$x^2=A+Bx$, $z=C+Dx$ と書くと $xz=0$ より $DA=0$, $C=-DB$ で、
$z^2=0$ の $x$ 係数より $D^2B=0$ である。従って
\[
    f=\operatorname{ord}D,\qquad e=\min(\operatorname{ord}A,\operatorname{ord}B),
    \qquad s\le 2f+e
\]
となる（位数は $s$ で打ち切る）。長さ $f$ では $z=0$ より鉛筆は $D$ 内にある。
長さ $e$ では $x^2=0$ なので固定接平面にある。実際前節までの局所座標で階数一行列は
定数行列、$x,z$ の二つの接行列、および $x^2$ の一行列の線型結合である。
いずれも二点因子全体が線型部分空間に入ればその相対線型包も入る。

こうして各因子の長さは、孤立点の固定平面、接平面が与える直線、および
$D$ の固定左核の直線での長さを順に重み $1,1,2$ で数えたもの以下である。
最後の直線が印以外に捉える長さは2以下であるから
\[
              13\le \operatorname{length} Z\le 6+3+2\cdot 2=13 .
\]
使った全ての不等式は等号である。固定左核に割り当てた $f$ の和はちょうど2で、
$f>0$ の各因子の長さは $s=2f+e\ge2f$ である（単根分離では $e=0$）。
よってその直線が捉える長さ2は非単純因子だけにあり、一つなら元の因子の長さは4以上である。
右核も同じ計算で扱える。
\end{proof}

\begin{lemma}\label{lem:confluent-three-first-jets}
十七点支持で三つの二重点があるとする。単純印の五空間の有限な階数一切断は、
対称性がないなら六つの被約点に限られる。
\end{lemma}
\begin{proof}
$J$ の外に十点以上ある。接次元2の節点と、曲線状で支持四点以下の切断は
それぞれ上界7と8で除かれる。五節点なら前節の $v,u_1,\ldots,u_4$ の模型である。
どちらの超平面も四節点以上を含めば単独の長さ4判定で対称性が出る。
内部節点が $1\le k\le3$ 個なら各内部節点から外に出る $5-k$ 本の割線の超平面交点は同じで、
外のグラフ点の総数は $11-k(4-k)\le8$ に落ちる。従って両超平面とも節点を避ける。
残余 $Z$ は各割線上の $B$ の剰余値で区別され、支持数は10または11である。
長さ2の因子は全て、相対割線が動く混合辺 $i$ にあり、元の二重点全体を使って $h(A)$ は単元である。

前節の十基底 $X_i,X_{ij}$、接線平方 $W_a^2$ と一次方向 $D_i=2x_i W_a$ を使う。
二本の交点は\emph{同じ}混合割線の変形から来る。すなわち $v+\epsilon v'$ を動かす速度は
両側で同じである。このため平方の同型が $X_i$ を $s_i X'_i$ に送れば、
$D_i$ は $s_iD'_i$ と $X'_i$ の倍の和に移る。これは両側で $a(v)=b(v)=1$ とした
交点式を同じ一様化変数で微分し、二重点上の単元を微分して得る。
$r_i=b_i/a_i$ と書く。$D_i$ の $X_j,X_{ij}$ 係数は $j\ne i$ で $a_j,-a_j$、
$i$ に接しない二添字成分は零である。

十基底が全て使われていれば二つ以上の可動混合辺 $i,j$ があるので、
それらの微分を比較すると $s_j=s_i r_j=s_{ij}=s_j r_i=s_i$。
よって $r_i=1$ であり、混合辺 $i$ の二交点が一致して $h(A)$ が単元であることに反する。
これは接線平方だけが欠ける場合も含む。他の場合は支持十点、三つの可動混合辺が必要であり、
全二重点をそこで使うので $Z$ の全ての支持で $h(A)$ は単元である。
欠けた基底 $X_m$ だけは $\tau X'_m+\sum_{j\ne m}\delta_j X'_j$ に送られ、
$\tau\ne0$ とする（ここで $j$ は十基底の添字）。

$m$ が単一添字 $k$ なら残りの三混合辺が全て可動である。
$D_i$ の $m$ 係数から $\tau a_k=s_i b_k$ が全ての可動 $i$ で成立し、$s_i$ は共通である。
可動二添字 $i,j$ に対して、第三の可動添字の微分の $ij$ 係数は両側で零なので
$\delta_{ij}=0$。従って $s_{ij}=s_i r_j=s_j r_i$ で $r_i=r_j$。
使われた単純節点間の割線 $ij$ で二超平面の交点が一致して矛盾する。

$m=12$ とする。$3,4$ がともに可動なら両者の微分は $m$ 係数が零なので、
先の十基底での二辺比較をそのまま適用できる。残るのは $1,2,3$ が可動の場合である。
$D_2$ の $13$ 係数は零で $m$ 係数は非零だから $\delta_{13}=0$。
従って $D_1,D_3$ の $13$ 係数から $s_{13}=s_1 r_3=s_3 r_1$。
一方 $D_3$ は $m$ 係数が零なので、その $1$ 係数によって $s_1=s_3 r_1$。
再び $r_3=1$ が単元条件に反する。
\end{proof}

\begin{theorem}\label{thm:finite-seventeen-site-pencils}
全中央支持4の有限窓で十七の幾何点を持つ曲線状点スキームについて、
反対階数10なら少なくとも片側対称であり、反対階数11は不可能である。
\end{theorem}
\begin{proof}
長さの分割は一つの四重点、三重点と二重点、または三つの二重点である。
単純印の五空間で対称性がなければ、正次元切断は先の接触次数補題の場合だけが残る。
実際、外向き接点だけの配置は $2+3+3$ 以下、孤立点が一つだけで接点のない場合は
$2+6$ 以下で、$J$ の外の十点に足りない。他の配置は
命題\ref{prop:curvilinear-marked-pencils}と定理\ref{thm:finite-single-triple-five-space}の
長さ4切断の議論による。
有限切断は被約六点または単純四点と二重節点の場合に限り、両超平面とも節点を避ける。
後者の証明は直前の $11-k(4-k)$ の計数をそのまま使える。
前者で十三本の割線による節点回避は曲線状窓すべてに使える。
長さ3以上のグラフ点ではどちらの場合も二行列の剰余値が比例する。
剰余で独立なら相対割線の変形と超平面の交点が閉埋込みでその局所環を受け取る必要があるが、
六節点では動かず、五節点でも混合割線の一階変形までしか動かないからである。

四重点だけなら接触次数補題の長さ2はその接線、三重点と二重点なら両者を結ぶ割線である。
各側においてこの固定直線には高々一つの単純印しか加えられない。
階数10または11なら各印で次元5以下の側を選ぶ。共通四空間なら対称性が得られ、
それ以外でも十五以上の単純印から高々二印を除けば、少なくとも七印が同じ側で有限切断を与える。
長さ3以上の点に印を移せば、その点自身を含め八つの異なる交点があり、共通四空間が強制される。

三つの二重点の場合、十四の単純印のそれぞれで次元5の側を一つ選ぶ。
各側の埋込み接点の例外印は二重点の三対の各割線で高々一印、従って高々三印である。
残る有限切断の印では補題\ref{lem:confluent-three-first-jets}により被約六節点となり、
三つの二重点\emph{全て}で二行列が比例する。
一つの側でこうした有限切断印が四つ以上あるとする。三つの二重点それぞれに印を取り直すと、
その二つの四空間の交叉を一方の変数の射影三空間に引き戻した線型空間には
同じ四つ以上の単純印が入り、さらに取り直した二重点自身が入る。
共通四空間なら対称性が得られるので、この線型空間は平面以下である。
四つの異なる印は共線になれないから平面を一意に張る。
従って三つの二重点は\emph{同じ}平面に全て入る。
この平面の長さ上界7より当該側の有限切断印は全体で四つ以下である。

各側で有限切断が高々四印、例外が高々三印なので十四印を覆うには
両側とも四印と三印で等号でなければならない。一方の側で得た平面は既に
四つの単純印と三つの二重点の剰余を含む。他方の例外印はこの位置で
二つの二重点を通る直線にある（接触次数補題は両核位置に成立する）。
この別の単純印も同じ平面に入り、長さ7を超える。これで階数11を含めて矛盾する。
\end{proof}
